\subsection{Canonical spectral Einstein functional and geometric Hodge closure}

The local target action is written as the low-curvature expansion of a canonical spectral
Einstein functional:
\begin{equation}
\label{eq:chi-effective-action}
\begin{aligned}
\Gammaeff
=
\int d^4x\sqrt{-g}\,
\Big[
c_R\chigeo^2R
+Z_\chi(\partial\chigeo)^2
-\mu_\Phi^2(\chigeo)\,\Phi^\dagger\Phi
+\lambda_\Phi(\Phi^\dagger\Phi)^2
+\frac{1}{6M^2}R^2
+\cdots
\Big].
\end{aligned}
\end{equation}
The relative bosonic endomorphism used below is
\begin{equation}
\label{eq:relative-bosonic-endomorphism}
E_{\mathrm{rel}}
=
-\frac14 R\,\mathbf 1_{48}
-
(\Phi^\dagger\Phi)\,P_2
+
\frac12\gamma^{\mu\nu}\mathbb F_{\mu\nu},
\end{equation}
and the resulting low-curvature residue decomposition is summarized in
\Cref{tab:heat-kernel-map}.

\begin{table}[t]
\centering
\scriptsize
\setlength{\tabcolsep}{3pt}
\begin{tabular}{@{}>{\raggedright\arraybackslash}p{0.15\linewidth}>{\raggedright\arraybackslash}p{0.22\linewidth}>{\raggedright\arraybackslash}p{0.22\linewidth}>{\raggedright\arraybackslash}p{0.14\linewidth}>{\raggedright\arraybackslash}p{0.14\linewidth}@{}}
\toprule
Residue term & Operator sector & Generated term & $\chigeo$ scaling & Role \\
\midrule
$a_0$ & relative bulk trace & vacuum-energy bookkeeping & $\chigeo^4$ & Hard \\
$a_2$ & seam-even bosonic bulk term & $\chigeo^2R$, $\chigeo^2\Phi^\dagger\Phi$ & $\chigeo^2$ & Closure \\
$a_4$ & local bosonic sector & $|D\Phi|^2$, $(\Phi^\dagger\Phi)^2$, $R\,\Phi^\dagger\Phi$, $\tr\mathbb F^2$, $R^2$ & $\chigeo^0$ & Closure \\
boundary/defect term & seam correction and parity data & local defect counterterms, seam-even selection & $\Delta_{\Seam}$ & Closure \\
$\eta$ term & APS spectral asymmetry & sign and spectral-flow control & scale free & Closure \\
\bottomrule
\end{tabular}
\caption{Low-curvature residue map for the canonical spectral Einstein functional.}
\label{tab:heat-kernel-map}
\end{table}

The resulting geometric logic is displayed in \Cref{fig:chi-bifurcation}. The canonical spectral
residue fixes the Einstein term first; the low-curvature expansion then yields the Higgs
normalization, the Higgs vacuum relation, and the retained curvature/gauge side terms. This is
why the local spectral-scale sector is the principal geometric closure channel of the manuscript
rather than just one more scalar parameter.

\begin{figure}[t]
\centering
\scriptsize
\begin{tikzpicture}[
    >=Latex,
    node distance=8mm and 11mm,
    a2/.style={draw=blue!60!black, fill=blue!6, rounded corners, align=center, text width=2.4cm, minimum height=0.9cm},
    a4/.style={draw=purple!70!black, fill=purple!7, rounded corners, align=center, text width=2.4cm, minimum height=0.9cm},
    flow/.style={draw=orange!75!black, fill=orange!8, rounded corners, align=center, text width=2.35cm, minimum height=0.95cm},
    outbox/.style={draw=green!50!black, fill=green!8, rounded corners, align=center, text width=2.6cm, minimum height=0.95cm}
]
\node[a2] (a2node) {canonical spectral residue\\Einstein seed};
\node[a4, right=24mm of a2node] (a4node) {low-curvature expansion\\Higgs / curvature terms};

\node[flow, below=of a2node] (ein0) {$\chigeo^2R$\\Einstein term first};
\node[flow, below=of a4node] (kin0) {$|D\Phi|^2,\ (\Phi^\dagger\Phi)^2$\\$\tr\mathbb F^2,\ R^2$};
\node[flow, below=of $(ein0)!0.5!(kin0)$, yshift=-3mm] (chi) {same seam-even $\chigeo$ channel\\canonical spectral functional};

\node[outbox, below left=of chi] (mpl) {$\Mpl^2=\chigeozero^2/(2\pi^2)$\\Einstein branch};
\node[outbox, below=of chi] (zh) {$Z_H=1/(16\pi^2)$\\canonical Higgs normalization};
\node[outbox, below right=of chi] (vev) {$\vgeo^2=Z_H\,\mu_\Phi^2(\chigeozero)/\lambda_\Phi$\\Higgs vacuum branch};
\node[flow, right=18mm of vev] (curv) {$R^2/M^2,\ \tr\mathbb F^2$\\higher-curvature / gauge branch};

\draw[->, thick] (a2node) -- (ein0);
\draw[->, thick] (a4node) -- (kin0);
\draw[->, thick] (ein0) -- (chi);
\draw[->, thick] (kin0) -- (chi);
\draw[->, thick] (chi) -- (mpl);
\draw[->, thick] (chi) -- (zh);
\draw[->, thick] (chi) -- (vev);
\draw[->, thick] (kin0.east) -- (curv.west);
\draw[->, thick] (zh) -- (vev);
\end{tikzpicture}
\caption{Canonical spectral flow of the local spectral-scale channel. The spectral residue fixes the
Einstein term first; the low-curvature expansion then yields Higgs normalization, the Higgs
vacuum branch, and the retained $R^2$/gauge side terms.}
\label{fig:chi-bifurcation}
\end{figure}

\begin{theorem}[Canonical spectral Einstein functional and scalaron seed]
\label{prop:canonical-bosonic-normalization}
\label{thm:canonical-einstein-residue}
In four dimensions the seam-even relative residue expansion gives
\[
\mathcal E_D
=
-\int d^4x\sqrt g
\left(
\frac{1}{24\pi^2}R
+
\frac{1}{48\pi^2}\Phi^\dagger\Phi
\right)
+
\mathcal E_\Sigma,
\]
and
\[
\begin{aligned}
\Gamma_D^{(4)}
={}&
\int d^4x\sqrt g
\Big(
\frac{1}{16\pi^2}|D\Phi|^2
+\frac{1}{16\pi^2}(\Phi^\dagger\Phi)^2
+\frac{1}{96\pi^2}R\,\Phi^\dagger\Phi
\\
&\qquad
+\frac{1}{192\pi^2}\tr\mathbb F^2
+\frac{1}{6M^2}R^2
+\cdots
\Big)
+\mathcal A_\Sigma.
\end{aligned}
\]
The seam-even scale channel $\chigeo$ is the unique scalar geometric mode on the retained branch;
together with the retained $R^2$ term it forms a positive scalaron block, so no independent
conformal negative mode survives as a separate physical degree of freedom.
Therefore
\[
\Gamma_{\mathrm{grav}}
:=
-6\chigeo^2\mathcal E_D+\Gamma_D^{(4)}
\]
has low-curvature form
\[
\Gamma_{\mathrm{grav}}
=
\int d^4x\sqrt g
\left(
\frac{\Mpl^2}{2}R
+
\frac{\Mpl^2}{12M^2}R^2
+
|D\Phi|^2
-V(\Phi,\chigeo)
-\sum_a \frac{1}{4g_a^2}\tr F_a^2
+\cdots
\right),
\]
with
\[
\Mpl^2=\frac{\chigeozero^2}{2\pi^2},
\qquad
Z_H=\frac{1}{16\pi^2},
\qquad
\Phi_{\mathrm{can}}:=\sqrt{Z_H}\,\Phi=\frac{\Phi}{4\pi},
\]
and therefore
\begin{equation}
\label{eq:induced-einstein}
\Mpl^2=\frac{\chigeozero^2}{2\pi^2},
\qquad
\vgeo^2=Z_H\,\frac{\mu_\Phi^2(\chigeozero)}{\lambda_\Phi}.
\end{equation}
After \Cref{thm:absolute-spectral-planck-closure}, the first identity is equivalently
\[
\frac{\Mpl^2}{\lambda_\Sigma^2}=\frac{\rho_\star}{2\pi^2}.
\]
\end{theorem}

\begin{proof}[Proof sketch]
The operator \Cref{eq:relative-bosonic-endomorphism} controls the seam-even residue expansion of
\Cref{def:canonical-residues}. Mellin expansion of the relative zeta pair fixes the order-one and
order-zero residues without any profile choice, so the Einstein term appears before the Higgs
kinetic, quartic, mixed-curvature, gauge-curvature, and $R^2$ terms. The coefficient of
$|D\Phi|^2$ is therefore the wave-function factor $Z_H=1/(16\pi^2)$. After the canonical rescaling
$\Phi_{\mathrm{can}}=\sqrt{Z_H}\,\Phi$, the Einstein and Higgs branches are read from one common
profile-free residue functional rather than from separate normalization conventions. Because the
same seam-even spectral channel carries $\chigeo$ and the retained $R^2$ block, the dangerous
independent conformal mode is absorbed into one positive scalaron sector.
\end{proof}

\begin{remark}[Boundary-spectral form of Newton's constant]
Using the standard reduced-Planck definition $\Mpl^{-2}=8\pi G_N$, the first relation in \Cref{eq:induced-einstein} is equivalently
\[
G_N\chigeozero^2=\frac{\pi}{4}.
\]
This Newton-form equation is therefore a rewriting of the Einstein-branch closure, not a second
independent prediction. In the boundary-normalized language of
\Cref{sec:absolute-dimensionless-metrology} it becomes
\[
G_N\lambda_\Sigma^2=\frac{\pi}{4\rho_\star},
\]
so Newton's constant is read as a boundary-spectral ratio rather than as an isolated SI datum.
\end{remark}

\begin{lemma}[Two-sheet seam-even zero-mode normalization]
\label{lem:two-sheet-zero-mode-normalization}
If the seam-even Higgs zero mode is normalized across both sheets of the oriented double cover, then its branch weight is
\[
w_H
=
2\,\gcar\,\beta_{\mathrm{rad}}^2
\exp\!\left[-\frac{\alpha^{-1}(0)+\delta_{\mathrm{ph}}}{5}\right].
\]
With
\[
\mu_\Phi^2(\chigeozero)=\frac{\chigeozero^2}{8\pi^2}w_H^2,
\qquad
\lambda_\Phi=\frac{1}{16\pi^2},
\]
the canonical Higgs vacuum expectation value is
\begin{equation}
\label{eq:carrier-dressed-v}
\vgeo
=
\Mpl\,\gcar\,\beta_{\mathrm{rad}}^2\,
\exp\!\left[-\frac{\alpha^{-1}(0)+\delta_{\mathrm{ph}}}{5}\right].
\end{equation}
\end{lemma}

\begin{proof}[Proof sketch]
The two-sheet seam-even normalization doubles the unnormalized zero-mode weight and therefore produces the stated $w_H$. Since the canonical Higgs field is $\Phi_{\mathrm{can}}=\sqrt{Z_H}\,\Phi$, its vacuum relation is
\[
\vgeo^2=\frac{Z_H\mu_\Phi^2(\chigeozero)}{\lambda_\Phi}
=
\frac{\chigeozero^2}{8\pi^2}w_H^2.
\]
Using $\Mpl^2=\chigeozero^2/(2\pi^2)$ from \Cref{prop:canonical-bosonic-normalization} gives $\vgeo=\Mpl w_H/2$, hence \Cref{eq:carrier-dressed-v}.
\end{proof}

\begin{corollary}[Hierarchy compression on the exact geometric branch]
\label{cor:hierarchy-compression-vgeo}
On the exact geometric branch one has
\[
\log\frac{\vgeo}{\Mpl}
=
\log\gcar
+
2\log\beta_{\mathrm{rad}}
-
\frac{\alpha^{-1}(0)+\delta_{\mathrm{ph}}}{5}.
\]
Hence the dominant suppression of the geometric Higgs source scale is carried by the electromagnetic
fixed point through the factor $e^{-\alpha^{-1}(0)/5}$. The theorem-level output is therefore the
dimensionless ratio $\vgeo/\Mpl$, equivalently the boundary-normalized ratio $\vgeo/\lambda_\Sigma$
or the hierarchy row $G_N\vgeo^2$, rather than a unitful GeV representative.
\end{corollary}

\begin{proof}
Taking the logarithm of \Cref{eq:carrier-dressed-v} gives the displayed identity immediately. The
dominant suppression is the exponential factor containing $\alpha^{-1}(0)$; the remaining terms are
the closed-branch carrier and retained-seed prefactors.
\end{proof}

\begin{remark}[Geometric versus benchmark electroweak scales]
For source formulas driven by the bosonic closure we write $\vgeo$ for the quantity in
\Cref{eq:carrier-dressed-v}. The preceding corollary concerns the internal geometric source scale
only and not yet the physical low-energy electroweak vacuum $v_{\mathrm{phys}}$. This is the
internal geometric branch value used in the UV source formulas and in the bosonic closure. The
physical electroweak vacuum parameter is not identified with $\vgeo$ itself; it is read only after
the renormalized low-energy branch has been formed. When the manuscript compares to a conventional
electroweak reference at the appendix interface, it uses a declared scheme-facing representative
$\vewref$. These three objects therefore belong to different
bookkeeping layers: $\vgeo$ to the internal geometric closure, $v_{\mathrm{phys}}$ to the physical
observable layer, and $\vewref$ only to the appendix comparison interface. No GeV representative is
part of the theorem-level gravitational or electroweak layer.
\end{remark}

\begin{definition}[Charged-current residue and physical electroweak scale]
\label{def:physical-ew-scale}
On the renormalized branch $\GrenTFPT$, let $\mathcal R_{\mathrm{CC}}^{\mathrm{TFPT}}(0)$ denote the
coefficient of the local charged-current four-fermion operator at zero momentum transfer, so that
\[
\Gamma_{\mathrm{TFPT}}^{\mathrm{ren}}
\supset
-\mathcal R_{\mathrm{CC}}^{\mathrm{TFPT}}(0)
\int d^4x\,
\bigl(\bar\nu_\mu\gamma^\mu P_L\mu\bigr)\bigl(\bar e\gamma_\mu P_L\nu_e\bigr).
\]
Define the physical electroweak scale by
\[
\mathcal R_{\mathrm{CC}}^{\mathrm{TFPT}}(0)
:=
\frac{1}{2v_{\mathrm{phys}}^2}
=
\frac{G_F^{\mathrm{TFPT}}}{\sqrt2}.
\]
Thus $\vgeo$ remains the geometric Higgs-closure scale, while $v_{\mathrm{phys}}$ is the low-energy
charged-current readout carried by $\GrenTFPT$.
\end{definition}

\begin{definition}[Branch source masses]
\label{def:branch-source-masses}
For each fermion sector define the source mass matrix
\[
M_f^{(0)}:=\frac{\vgeo}{\sqrt2}\,Y_f^\star.
\]
These are branch masses on the exact $\vgeo$ layer and not yet physical pole masses.
\end{definition}

\begin{theorem}[Native electroweak source-to-physical matching]
\label{thm:ew-geometric-to-physical-matching}
Let the renormalized charged-current residue at zero momentum transfer be written as
\[
\mathcal R_{\mathrm{CC}}^{\mathrm{TFPT}}(0)
=
\frac{g_2^2}{8m_W^2}\,(1+\Delta r_{\mathrm{CC}}^{\mathrm{TFPT}}).
\]
Write the renormalized $W$ mass in the form
\[
m_W^2=
\frac{g_2^2}{4}Z_\Phi(0)(\vgeo+\delta v)^2\,(1+\delta_W).
\]
Define
\[
\mathcal Z_{\mathrm{EW}}^{\mathrm{TFPT}}
:=
 \frac{1}{2\vgeo^2\,\mathcal R_{\mathrm{CC}}^{\mathrm{TFPT}}(0)}
 =
Z_\Phi(0)\left(1+\frac{\delta v}{\vgeo}\right)^2
\,\frac{1+\delta_W}{1+\Delta r_{\mathrm{CC}}^{\mathrm{TFPT}}}.
\]
Then the physical electroweak scale satisfies
\[
v_{\mathrm{phys}}
=
\vgeo\sqrt{\mathcal Z_{\mathrm{EW}}^{\mathrm{TFPT}}},
\qquad
v_{\mathrm{phys}}^2
=
Z_\Phi(0)(\vgeo+\delta v)^2\,\frac{1+\delta_W}{1+\Delta r_{\mathrm{CC}}^{\mathrm{TFPT}}}.
\]
Equivalently,
\[
\mathcal Z_{\mathrm{EW}}^{\mathrm{TFPT}}
=
\mathcal R_{\mathrm{CC}}\!\left[\Gamma_{\mathrm{TFPT}}^{\mathrm{ren}}\right].
\]
The charged-current residue and the four factors entering $\mathcal Z_{\mathrm{EW}}^{\mathrm{TFPT}}$
are canonically determined by the renormalized 1PI package $\Gamma_{\mathrm{TFPT}}^{\mathrm{ren}}$;
hence
$\mathcal Z_{\mathrm{EW}}^{\mathrm{TFPT}}$ is a branch invariant and not a fit parameter.
\end{theorem}

\begin{proof}
From
\[
\mathcal R_{\mathrm{CC}}^{\mathrm{TFPT}}(0)=\frac{1}{2v_{\mathrm{phys}}^2}
\]
and
\[
\mathcal R_{\mathrm{CC}}^{\mathrm{TFPT}}(0)=\frac{g_2^2}{8m_W^2}(1+\Delta r_{\mathrm{CC}}^{\mathrm{TFPT}})
\]
one gets
\[
\frac{1}{2v_{\mathrm{phys}}^2}
=
\frac{g_2^2}{8m_W^2}(1+\Delta r_{\mathrm{CC}}^{\mathrm{TFPT}}).
\]
Substituting
\[
m_W^2=
\frac{g_2^2}{4}Z_\Phi(0)(\vgeo+\delta v)^2(1+\delta_W)
\]
yields
\[
\frac{1}{2v_{\mathrm{phys}}^2}
=
\frac{1+\Delta r_{\mathrm{CC}}^{\mathrm{TFPT}}}
{2Z_\Phi(0)(\vgeo+\delta v)^2(1+\delta_W)}.
\]
Solving for $v_{\mathrm{phys}}$ gives the displayed identity. Rewriting the result as
\[
\mathcal Z_{\mathrm{EW}}^{\mathrm{TFPT}}
=
\bigl(2\vgeo^2\mathcal R_{\mathrm{CC}}^{\mathrm{TFPT}}(0)\bigr)^{-1},
\]
one obtains $v_{\mathrm{phys}}=\vgeo\sqrt{\mathcal Z_{\mathrm{EW}}^{\mathrm{TFPT}}}$. By
\Cref{thm:operational-completeness}, the charged-current residue, the Higgs wave-function factor
$Z_\Phi(0)$, the vacuum shift $\delta v$, the $W$-mass correction $\delta_W$, and the charged-current
remainder $\Delta r_{\mathrm{CC}}^{\mathrm{TFPT}}$ are all read from the same renormalized
admissible 1PI action $\Gamma_{\mathrm{TFPT}}^{\mathrm{ren}}$, so no additional continuous fit
parameter enters.
\end{proof}
